\documentclass[11pt,a4paper]{article}
\usepackage[T1]{fontenc}
\usepackage{lmodern}
\usepackage[margin=26mm]{geometry}
\usepackage{amsmath,amssymb,amsthm,booktabs,array,longtable}
\usepackage{microtype,multicol}
\usepackage[hidelinks]{hyperref}
\setlength{\parskip}{3pt}
\setlength{\emergencystretch}{2em}
\raggedbottom
\allowdisplaybreaks[1]
\clubpenalty=10000
\widowpenalty=10000
\newtheorem{theorem}{Theorem}[section]
\newtheorem{lemma}[theorem]{Lemma}
\newtheorem{corollary}[theorem]{Corollary}
\theoremstyle{remark}
\newtheorem{remark}[theorem]{Remark}
\newcommand{\Z}{\mathbb Z}
\newcommand{\Q}{\mathbb Q}
\newcommand{\N}{\mathbb N}
\newcommand{\F}{\mathbb F}
\newcommand{\ps}[1]{#1[[x]]}
\newcommand{\Rev}{\operatorname{Rev}}
\newcommand{\rev}{^{\langle-1\rangle}}
\newcommand{\Plus}{\operatorname{Plus}}
\newcommand{\Minus}{\operatorname{Minus}}
\newcommand{\coeff}[2]{[x^{#1}]#2}
\newcommand{\E}{\mathcal E}
\Urlmuskip=0mu plus 1mu
\hypersetup{pdftitle={A389537: Residues of an alternating lacunary inverse},pdfauthor={Dev Goyal},pdfsubject={Separate manuscript with AI disclosure and Lean proof map}}
\title{A389537\\[4pt]\large Residues of an alternating lacunary inverse}
\author{Dev Goyal}
\date{24 October 2026}
\begin{document}
\maketitle
\begin{center}
\fbox{\begin{minipage}{0.9\linewidth}\small
\textbf{AI disclosure.} ChatGPT (OpenAI) generated the mathematical arguments,
Lean proofs, and this manuscript in response to Dev Goyal's requests.
Dev Goyal is the sole human author. The source sequence and conjecture or
observation are credited to Paul D. Hanna.
\textbf{Independent human review is pending.}
\end{minipage}}
\end{center}
\begin{abstract}
The unique normalized integer solution of $A(A-x)=A^2$ has residue two modulo four exactly at indices $3\cdot2^k$ above degree two. We identify it as a genuine compositional inverse and prove the complete pattern by exact integral halving followed by a characteristic-two recursion.
\end{abstract}

\section{Source, notation, and scope}

The source is Paul D. Hanna's OEIS entry A389537 \cite{oeis537},
retrieved on 24 October 2026. Its initial sequence index is 1.
We also retain the formal zero constant coefficient at index zero.

All series are formal. The set $\N$ includes zero, and $[x^n]U$ denotes a
coefficient. A normalized series is $x+O(x^2)$; $O(x^r)$ means divisibility
by $x^r$, over any coefficient ring used below. We write $U\circ V=U(V(x))$
only when $V(0)=0$. Compositional inversion, denoted $\Rev(U)$ or $U\rev$,
is distinct from multiplicative inversion of a unit. Congruences are
coefficientwise; ordinary products and compositional iterates are distinct.

\section{Formal algebra used in the proof}

\begin{lemma}[Coefficientwise construction]\label{lem:contraction}
Let $K$ be a commutative ring. Suppose an operator $\Phi$ preserves
normalized series and, for $r\ge2$, satisfies
\[
 G-H=O(x^r)\quad\Longrightarrow\quad
 \Phi(G)-\Phi(H)=O(x^{r+1}).
\]
Then $\Phi$ has exactly one normalized fixed point in $\ps K$.
More generally, a residual equation whose degree-$r$ condition determines
the next unknown coefficient with multiplier one has a unique solution
once its initial coefficients are fixed and its lower-degree conditions
hold identically.
\end{lemma}

\begin{proof}
Begin with $G_0=x$ and set $G_{m+1}=\Phi(G_m)$. Normalization gives
$G_1-G_0=O(x^2)$. Inductively,
$G_{m+1}-G_m=O(x^{m+2})$, so each coefficient eventually stabilizes.
Define $G$ using those stable coefficients. The displayed implication
ensures that each coefficient of $\Phi(G)$ equals the corresponding stable
coefficient, hence $G=\Phi(G)$. If two fixed points agree below degree
$r$, they agree through degree $r$; induction proves uniqueness.

For the residual version, set the next coefficient to the negative of the
known residual term. Its multiplier is one, so the construction uses only
ring operations. No earlier equation changes, and every equation is
eventually imposed. The same coefficient comparison proves uniqueness.
\end{proof}

\begin{lemma}[Units and reversion]\label{lem:inverse}
Every constant-one series has a unique constant-one multiplicative
inverse. Every normalized series has a unique normalized two-sided
compositional inverse. These inversions, formal differentiation, and
valid substitution commute with coefficient ring homomorphisms.
\end{lemma}

\begin{proof}
For $UV=1$ with $U(0)=1$, the coefficient of degree $n>0$ is $v_n$ plus
a finite expression in earlier coefficients of $V$, which determines
$v_n$ without division. For normalized $G$, the coefficient of degree
$n>1$ in $G(H)$ is $h_n$ plus an expression in earlier coefficients of
$H$, giving a unique normalized right inverse $H$.

To see the second inverse identity, construct a normalized right inverse
$J$ of $H$. Associativity gives
$G=G\circ(H\circ J)=(G\circ H)\circ J=J$, and hence $H\circ G=x$.
The same association proves uniqueness. Each coefficient of a valid
composition is a finite polynomial in coefficients. Differentiation
multiplies the degree-$(n+1)$ coefficient by $n+1$. Ring homomorphisms
therefore preserve these operations and equations; uniqueness identifies
the mapped inverses with the inverses of the mapped series.
\end{proof}

\begin{lemma}[Binary support]\label{lem:binary}
There is a unique normalized integral series $C$ satisfying $C=x+C^2$.
Its reduction modulo two is
\[
 P(x)=\sum_{k\ge0}x^{2^k}.
\]
More generally, for every integer $d\ge1$, the unique zero-constant
solution of $V+V^2=x^d$ in $\ps{\F_2}$ is
\[
 V=P(x^d)=\sum_{k\ge0}x^{d2^k}.
\]
\end{lemma}

\begin{proof}
The operator $G\mapsto x+G^2$ satisfies Lemma~\ref{lem:contraction},
since $G^2-H^2=(G-H)(G+H)$ and $G+H=O(x)$. In characteristic two,
every series $V$ satisfies $V(x)^2=V(x^2)$. This follows by truncating
above any fixed degree and applying the polynomial Frobenius identity.

Thus $V+V^2=x^d$ is a recursion in which the coefficient at an odd positive
index is fixed immediately, while the coefficient at an even index $n$
depends on the coefficient at $n/2<n$. The constant coefficient is fixed
to zero. This proves uniqueness. The proposed sum is well-defined
coefficientwise because its exponents tend to infinity; adding its square
cancels all terms except $x^d$. Set $d=1$ to identify the reduction of $C$.
\end{proof}

\begin{lemma}[Doubling and substitution]\label{lem:double}
For $U,V\in(\Z/4\Z)[[x]]$,
\[
 2U=2V\quad\Longleftrightarrow\quad
 U\equiv V\pmod2.
\]
If $B,D$ have zero constant coefficient and $B\equiv D\pmod2$, then
$2H(B)=2H(D)$ for every series $H$ over $\Z/4\Z$.
\end{lemma}

\begin{proof}
For residues $u,v$ modulo four, $2u=2v$ exactly when their reductions modulo
two agree. Apply this fact to each coefficient. Reduction modulo two
commutes with substitution, giving the second assertion. In particular,
this lemma never cancels two inside $\Z/4\Z$.
\end{proof}

\section{A389537: an alternating lacunary inverse}

\subsection{Construction and the original equation}

Over any commutative ring, define
\begin{equation}\label{eq:lacunaryR}
 R(x)=\sum_{k\ge0}(-1)^kx^{2^k}.
\end{equation}
This is a normalized formal series, and cancellation of successive terms
gives $R(x)+R(x^2)=x$. It is the unique normalized solution of this
equation: an odd degree above one has coefficient zero, and an even
degree is determined by minus the coefficient at half that degree.
Equivalently, $U\mapsto x-U(x^2)$ strictly increases the order of
differences between normalized inputs, so Lemma~\ref{lem:contraction}
applies. In Lean the latter construction is used.

Let $A=R\rev$ be the genuine normalized compositional inverse.
Substituting $A$ in $R+R(x^2)=x$ gives
\[
 x+R(A^2)=A.
\]
Apply $A$ to $R(A^2)=A-x$ to obtain the literal A389537 equation
\begin{equation}\label{eq:lacunarysource}
 A(A(x)-x)=A(x)^2.
\end{equation}
Conversely, let normalized $U$ satisfy \eqref{eq:lacunarysource},
and let $V=U\rev$. Applying $V$ on the outside gives
$U-x=V(U^2)$. Substituting $V$ for $x$ then gives
$x-V=V(x^2)$. The uniqueness of \eqref{eq:lacunaryR} implies $V=R$
and hence $U=A$. We have therefore constructed the unique normalized
integral solution of the original equation, not just an auxiliary
series with a related congruence. Coefficient maps preserve this
construction by Lemma~\ref{lem:inverse}.

\subsection{Parity, exact halving, and the modulus-four result}

\begin{theorem}\label{thm:lacunary}
For the preceding integer source series, $a_1=a_2=1$, and for every $n>2$,
\[
 a_n\equiv
 \begin{cases}
  2\pmod4,&n=3\cdot2^k\text{ for some }k\in\N,\\
  0\pmod4,&\text{otherwise}.
 \end{cases}
\]
\end{theorem}

\begin{proof}
Modulo two, Frobenius changes the inverse equation to $R+R^2=x$.
Thus $(x+x^2)\circ R=x$. The uniqueness of normalized compositional
inverses shows
\begin{equation}\label{eq:lacunaryparity}
 A\equiv x+x^2\pmod2.
\end{equation}
In particular, there is a well-defined integral series
\begin{equation}\label{eq:half}
 H=\frac{A-x-x^2}{2}\in\ps{\Z},\qquad H(0)=0.
\end{equation}
This division takes place coefficientwise in the integers, after the
divisibility proof \eqref{eq:lacunaryparity}; it is not division by two
inside $\Z/4\Z$. The initial values $a_1=a_2=1$ also follow directly
from $R=x-x^2+O(x^4)$ and $R(A)=x$.

Reduce \eqref{eq:half} modulo four, retaining the notation $H$, and set
$D=A-x=x^2+2H$. Both $D$ and $x^2$ have zero constant coefficient,
and they have the same reduction modulo two. Lemma~\ref{lem:double}
therefore gives $2H(D)=2H(x^2)$. Expanding the source equation yields
\begin{align*}
 A(D)&=D+D^2+2H(D)
       =x^2+2H+x^4+2H(x^2),\\
 A^2&=(x+x^2+2H)^2=x^2+2x^3+x^4.
\end{align*}
All equalities in this calculation are in $(\Z/4\Z)[[x]]$.
After subtraction,
\[
 2\bigl(H+H(x^2)\bigr)=2x^3.
\]
Applying the equivalence in Lemma~\ref{lem:double}, and writing
$\overline H$ for reduction modulo two, gives
\[
 \overline H+\overline H(x^2)=x^3,
 \qquad\text{equivalently}\qquad
 \overline H+\overline H^{\,2}=x^3.
\]
Since $\overline H(0)=0$, Lemma~\ref{lem:binary} identifies it uniquely:
\[
 \overline H=\sum_{k\ge0}x^{3\cdot2^k}.
\]
For $n>2$, equation \eqref{eq:half} gives $a_n\equiv2[x^n]H\pmod4$.
The displayed support proves both residue cases.
\end{proof}

This is the modulus-four conjecture in A389537 \cite{oeis537}.
The same entry's empirical observation about lengths of sign runs is a
separate problem and is not a consequence claimed here.


\section{Separate Lean proof and checked scope}

The entry file is \nolinkurl{lean/Entries/A389537.lean}. Its explicit
theorem statements and proof terms use the namespace
\nolinkurl{FinalCollection.A389537}. The complete imported implementations
and their shared mathematical dependencies are included in the collection.
The following local theorem names give the proof map:

\begin{multicols}{2}\raggedright\begin{itemize}\setlength{\itemsep}{2pt}\setlength{\parsep}{0pt}
\item \nolinkurl{auxiliary_exists_unique}.
\item \nolinkurl{auxiliary_equation}.
\item \nolinkurl{inverse_right}.
\item \nolinkurl{inverse_left}.
\item \nolinkurl{normalized}.
\item \nolinkurl{source_equation}.
\item \nolinkurl{source_exists_unique}.
\item \nolinkurl{source_unique}.
\item \nolinkurl{series_mod_two}.
\item \nolinkurl{integer_decomposition}.
\item \nolinkurl{half_equation_mod_two}.
\item \nolinkurl{half_binary_support}.
\item \nolinkurl{residue_mod_four}.
\end{itemize}\end{multicols}

The main implementation module is \nolinkurl{Round3.LacunaryInverse}.
The catalogue records the full original endpoint names and the new entry
theorems. The project pins Lean 4.19.0 and Mathlib revision
\nolinkurl{c44e0c8ee63ca166450922a373c7409c5d26b00b}. The collection
report records source recompilation and transitive axiom checks. Accepted
axioms are propositional extensionality, classical choice, and quotient
soundness. All index conclusions above are unbounded; exact finite
coefficient checks in the package provide separate corroboration.

\paragraph{Limits and status.}
The separate empirical sign-run observation is not proved. Integer halving precedes reduction modulo four.
The manuscript consolidates completed results; it does not assert global
novelty, exclusive priority, peer review, OEIS acceptance, or a new submission.
Standard algebraic ingredients are not claimed as new discoveries.

\begin{thebibliography}{9}
\raggedright
\bibitem{oeis537}
Paul D. Hanna, \emph{OEIS A389537}, 3 November 2025.
\url{https://oeis.org/A389537}. Retrieved 24 October 2026.
\end{thebibliography}
\end{document}
